\scanpage{43}
\mspage{4}
\uline{\unsure{Théorèmes|Théorème} de rigidité}

\margin{\uline{Déf}} Un préschéma en \unsure{groupe|groupes} commutatif \inserted{$G$} de type fini sur $\underline S$ est dit \uline{rigide} si \deleted{pour tt}

(i) $G$ est \deleted{plat} simple sur $S$

(ii) Pour tt $s\in S$, et \deleted{tout \unsure{entier} \unsure{n}}\inserted{nombre} premier $\ell$ \deleted{à la caractéristique} distinct de $\operatorname{car} k(s)$, $\overline G_s$ est \deleted{le plus petit sous-schéma fermé contenant le \unsure{$ {}_{\ell^k}(G_s)$}.}\deleted{\inserted{de $G_s$ \unsure{contenant}}} l'adhérence des $ {}_{\ell^k}\overline G_s$ [où $\overline G_s$ est obtenu par $G_s$ en passant à \deleted{\unsure{un}} $\overline{k(s)}$]

% slanted
\margin{Il suffit \unsure{de|ici} l'exiger pour \uline{un} $\ell$}
(iii) Pour tt \deleted{\unsure{$s\in S$, tt}} ouvert $U$ de $S$, tt \deleted{\unsure{entier}} $n$ premier aux \deleted{\unsure{car.}} \unsure{car.} des pts de $U$, le schéma $ {}_n(G|U)$ est \uline{fini} sur $S$\unsure{]}

\margin{\uline{Prop}} Remarquons qu'on peut prouver qu'il \deleted{faut et} suffit pour \deleted{\unsure{cela}}\inserted{(ii)} que \deleted{$G_s$ \unsure{ne contienne} \inserted{dans les} \unsure{deux}} \inserted{parmi} \deleted{\unsure{deux|dans}} les facteurs simples $G_a$, $G_m$, $Ab$, $\mathbb Z/\ell$, $\mathbb Z/p$ qui interviennent en $G_{\bar s}=G$, ne figurent ni $G_a$, ni $\mathbb Z/p$. \deleted{\unsure{cf deux}}

% slanted
\margin{Développer le sorite : extension de \unsure{grpes}, ext. de \unsure{fais} \unsure{homog.}, \deleted{\illegible{1}} --- fid. plats}
\deleted{En effet, on (et cette condition est suffisante)} \deleted{Pour la nécessité,} On \deleted{sait}\inserted{prouve pour cela} qu'une extension de groupes rigides sur $k$ alg clos est rigide ; puis \deleted{\unsure{q}} que $G_m$, $Ab$ sont rigides (par des démonstrations quasi-identiques), $\mathbb Z/\ell$ \unsure{l'étant} trivialement. \deleted{\unsure{cf}}\inserted{On} utilisera alors le th. de structure des groupes commutatifs \unsure{alg.} Pour \deleted{la nécessité,}\inserted{la réciproque,}

\margin{\diagram{$\begin{array}{ccc}G_m&&\mathbb Z_{\ell}\\&Ab&\\G_a&&\mathbb Z_p\end{array}$}}
% archive sheet number: 42
